% =============================================================================
% Guide d'implémentation — Mode hors ligne (Offline)
% Yowyob ERP — Comptabilité Générale et Analytique
% =============================================================================
\documentclass[12pt,a4paper]{report}

\usepackage[utf8]{inputenc}
\usepackage[T1]{fontenc}
\usepackage[english]{babel}
\usepackage{geometry}
\usepackage{longtable}
\usepackage{array}
\usepackage{fancyhdr}
\usepackage{amsmath}

\geometry{margin=2.5cm}

\pagestyle{fancy}
\fancyhf{}
\fancyhead[L]{\footnotesize Guide Offline — Yowyob ERP}
\fancyhead[R]{\footnotesize Backend \& Frontend}
\fancyfoot[C]{\thepage}
\renewcommand{\headrulewidth}{0.4pt}

\begin{document}

% -----------------------------------------------------------------------------
\begin{titlepage}
  \centering
  \vspace*{2cm}
  {\LARGE\textbf{YOWYOB ERP}\\[0.4em]Comptabilité Générale et Analytique\par}
  \vspace{1.8cm}
  {\Huge\textbf{Guide d'implémentation\\[0.35em]Mode hors ligne (Offline)}\par}
  \vspace{1.2cm}
  {\large Spécification technique Backend et Frontend\par}
  \vspace{2.2cm}
  \begin{tabular}{rl}
    \textbf{Version} & 1.0 \\
    \textbf{Date} & Juillet 2026 \\
    \textbf{Public} & Développeurs frontend, développeurs backend, architectes \\
    \textbf{Statut} & Document de référence pour la mise en œuvre \\
  \end{tabular}
  \vfill
  {\small Ce document décrit l'architecture, les contrats API et les procédures d'implémentation pour un fonctionnement \emph{online-first avec repli offline}.\par}
\end{titlepage}

\tableofcontents
\newpage

% =============================================================================
\chapter{Introduction}
% =============================================================================

\section{Objectif}

Ce guide définit comment permettre à un utilisateur de continuer à travailler dans l'application web lorsque la connexion réseau est indisponible, puis de retrouver automatiquement l'ensemble de son travail dans l'application une fois la connexion rétablie.

Le principe retenu est le suivant :

\begin{itemize}
  \item \textbf{En ligne} : communication systématique avec le backend (source de vérité serveur).
  \item \textbf{Hors ligne} : bascule transparente vers une mémoire locale persistante.
  \item \textbf{Retour en ligne} : synchronisation automatique de la file d'attente vers le backend.
\end{itemize}

\section{Périmètre}

\begin{longtable}{@{}p{4cm}p{9cm}@{}}
\hline
\textbf{Inclus} & Détection réseau, cache local, file d'attente (outbox), moteur de sync, contrat API backend, gestion des conflits, indicateur discret, pilote écritures analytiques \\
\hline
\textbf{Hors périmètre initial} & PWA complète (Service Worker), sync multi-appareils temps réel, offline sur pièces jointes volumineuses \\
\hline
\end{longtable}

\section{Principes directeurs}

\begin{enumerate}
  \item \textbf{Online-first} : tant que le réseau et l'API répondent, toute lecture et écriture passe par le backend.
  \item \textbf{Offline en dernier recours} : le stockage local n'est activé que si la connexion ou l'API est indisponible.
  \item \textbf{Transparence} : l'utilisateur n'est pas bloqué ; un indicateur discret signale les éléments en attente de synchronisation.
  \item \textbf{Pas de perte} : aucune donnée locale n'est supprimée avant accusé de réception serveur.
  \item \textbf{Idempotence} : le backend doit accepter le rejeu d'opérations sans créer de doublons.
\end{enumerate}

% =============================================================================
\chapter{Architecture globale}
% =============================================================================

\section{Vue d'ensemble}

\begin{verbatim}
┌─────────────────────────────────────────────────────────────────┐
│                        INTERFACE UTILISATEUR                     │
└───────────────────────────────┬─────────────────────────────────┘
                                │
                    ┌───────────▼───────────┐
                    │  Couche Réseau        │
                    │  (détection online)   │
                    └───────────┬───────────┘
                                │
              ┌─────────────────┴─────────────────┐
              │ EN LIGNE                        │ HORS LIGNE
              ▼                                 ▼
    ┌─────────────────┐               ┌─────────────────┐
    │  API Backend    │               │  IndexedDB      │
    │  (REST)         │               │  Cache + Outbox │
    └────────┬────────┘               └────────┬────────┘
             │                                  │
             └──────────────┬───────────────────┘
                            ▼
                 ┌─────────────────────┐
                 │  Moteur de sync     │
                 │  (flush outbox)     │
                 └─────────────────────┘
\end{verbatim}

\section{Composants frontend}

\begin{longtable}{@{}p{3.5cm}p{9.5cm}@{}}
\hline
\textbf{Composant} & \textbf{Rôle} \\
\hline

\texttt{network-status} & Détecte online/offline via \texttt{navigator.onLine}, événements navigateur et échecs API consécutifs \\
\texttt{offline/db} & Schéma IndexedDB (Dexie) : entités, outbox, métadonnées \\
\texttt{offline/outbox} & Enregistre chaque mutation (CREATE, UPDATE, DELETE) en attente \\
\texttt{offline/sync-engine} & Rejoue la file d'attente au retour en ligne \\
\texttt{offline/handlers} & Adaptateurs par type d'entité (écritures analytiques, etc.) \\
\texttt{use-network-status} & Hook React exposant \texttt{isOnline}, \texttt{pendingCount} \\
\texttt{OfflineProvider} & Provider global monté dans le layout dashboard \\
\hline
\end{longtable}

\section{Composants backend}

\begin{longtable}{@{}p{3.5cm}p{9.5cm}@{}}
\hline
\textbf{Composant} & \textbf{Rôle} \\
\hline

API REST CRUD & Endpoints standards par entité métier \\
Idempotence & Acceptation de \texttt{clientMutationId} ou en-tête \texttt{Idempotency-Key} \\
Versionnement & Champ \texttt{updatedAt} et/ou \texttt{version} sur chaque entité \\
Sync delta & Endpoint \texttt{GET ?since=timestamp} (phase 2) \\
Gestion conflits & Réponse HTTP 409 avec détail structuré \\
Audit & Journal des opérations comptables \\
\hline
\end{longtable}

% =============================================================================
\chapter{Détection de l'état réseau}
% =============================================================================

\section{Signaux utilisés}

\begin{enumerate}
  \item \texttt{navigator.onLine} : détection immédiate de coupure réseau locale.
  \item Événements \texttt{window.online} et \texttt{window.offline}.
  \item Compteur d'échecs API consécutifs (timeout, \texttt{Failed to fetch}, status 0).
  \item Repassage en mode online uniquement après un appel API réussi.
\end{enumerate}

\section{Distinction erreur réseau vs erreur métier}

\begin{longtable}{@{}p{3cm}p{4cm}p{6cm}@{}}
\hline
\textbf{Type} & \textbf{Exemples} & \textbf{Comportement} \\
\hline

Réseau & Timeout, fetch failed, status 0 & Bascule offline, cache + outbox \\
Métier & 422 validation, 403 interdit & Message utilisateur, pas de bascule offline \\
Serveur & 500, 502, 503 & Retry avec backoff, pas de bascule immédiate \\
Conflit & 409 & Marquer conflit, conserver donnée locale \\
\hline
\end{longtable}

\section{Algorithme de décision}

\begin{verbatim}
function shouldUseOffline(): boolean {
  if (!navigator.onLine) return true;
  if (consecutiveNetworkFailures >= 2) return true;
  return false;
}
\end{verbatim}

% =============================================================================
\chapter{Implémentation Frontend}
% =============================================================================

\section{Stockage local : IndexedDB}

\textbf{Ne pas utiliser} \texttt{localStorage} comme stockage principal offline (limite 5 Mo, pas de requêtes structurées).

\subsection{Schéma Dexie (version 1)}

\begin{verbatim}
entities  : key, entity, entityId, updatedAt
outbox    : id, entity, entityId, status, createdAt
meta      : key
\end{verbatim}

\subsection{Table \texttt{entities}}

Cache des données métier pour consultation hors ligne.

\begin{verbatim}
{
  key: "ecriture_analytique:ea-uuid",
  entity: "ecriture_analytique",
  entityId: "ea-uuid",
  data: { ... },
  updatedAt: "2026-07-08T12:00:00Z",
  syncStatus: "synced" | "pending" | "local_only"
}
\end{verbatim}

\subsection{Table \texttt{outbox}}

File d'attente des opérations à envoyer au backend.

\begin{verbatim}
{
  id: "op-uuid",
  entity: "ecriture_analytique",
  action: "CREATE" | "UPDATE" | "DELETE",
  entityId: "ea-uuid",
  payload: { ... },
  clientMutationId: "cm-uuid",
  createdAt: "2026-07-08T12:00:00Z",
  status: "pending" | "syncing" | "done" | "failed" | "conflict",
  retries: 0,
  lastError: null
}
\end{verbatim}

\section{Flux de lecture}

\begin{enumerate}
  \item Si \textbf{en ligne} : appel API $\rightarrow$ mise à jour cache IndexedDB $\rightarrow$ affichage.
  \item Si \textbf{hors ligne} : lecture cache IndexedDB $\rightarrow$ affichage.
  \item Si API échoue (réseau) : lecture cache en repli, bascule offline.
\end{enumerate}

\section{Flux d'écriture}

\begin{enumerate}
  \item Validation métier légère côté frontend (champs obligatoires, formats).
  \item Si \textbf{en ligne} : envoi API $\rightarrow$ si succès, mise à jour cache ; si échec réseau $\rightarrow$ étape 3.
  \item Si \textbf{hors ligne} ou échec réseau :
    \begin{itemize}
      \item Sauvegarde immédiate dans IndexedDB.
      \item Ajout dans outbox (status \texttt{pending}).
      \item Mise à jour UI (optimistic).
    \end{itemize}
  \item Au retour en ligne : \texttt{syncEngine.flush()} rejoue l'outbox dans l'ordre chronologique.
\end{enumerate}

\section{Règle d'or}

\begin{quote}
\textbf{Ne jamais supprimer une entrée de l'outbox ni une entité locale tant que le serveur n'a pas confirmé la réception (HTTP 200/201).}
\end{quote}

\section{Moteur de synchronisation}

\subsection{Algorithme flush}

\begin{enumerate}
  \item Récupérer les opérations \texttt{outbox} où \texttt{status = pending} ou \texttt{failed} (retries < 5), triées par \texttt{createdAt}.
  \item Pour chaque opération :
    \begin{itemize}
      \item Passer en \texttt{syncing}.
      \item Appeler le handler de l'entité (\texttt{push}).
      \item Si succès : \texttt{done}, mettre à jour \texttt{entities.syncStatus = synced}.
      \item Si erreur réseau : \texttt{pending}, incrémenter \texttt{retries}, arrêter le flush.
      \item Si 409 conflit : \texttt{conflict}, notifier discrètement.
      \item Si 422 métier : \texttt{failed}, conserver localement.
    \end{itemize}
  \item Émettre événement \texttt{sync:complete} avec le nombre d'opérations synchronisées.
\end{enumerate}

\subsection{Backoff}

\begin{longtable}{@{}cc@{}}
\hline
\textbf{Tentative} & \textbf{Délai avant retry} \\
\hline
1 & 0 s (immédiat au online) \\
2 & 2 s \\
3 & 5 s \\
4 & 15 s \\
5 & 30 s \\
\hline
\end{longtable}

\section{Intégration dans l'application existante}

\subsection{Layout dashboard}

Monter \texttt{OfflineProvider} qui :
\begin{itemize}
  \item initialise IndexedDB ;
  \item écoute online/offline ;
  \item lance le flush au retour en ligne ;
  \item expose le nombre d'opérations en attente.
\end{itemize}

\subsection{useAutoRefresh}

Désactiver le polling quand hors ligne :

\begin{verbatim}
useAutoRefresh(callback, deps, { enabled: isOnline });
\end{verbatim}

\subsection{Indicateur utilisateur}

Point discret dans le header :
\begin{itemize}
  \item Vert : en ligne, rien en attente.
  \item Orange : hors ligne OU opérations en attente de sync.
  \item Tooltip : « 3 élément(s) en attente de synchronisation ».
\end{itemize}

\section{Pilote : écritures analytiques}

Première entité migrée :
\begin{itemize}
  \item Remplacement de \texttt{localStorage} par IndexedDB.
  \item Migration automatique des données existantes au premier chargement.
  \item Chaque CREATE/UPDATE/DELETE alimente l'outbox.
  \item Handler sync prêt à brancher l'API backend quand disponible.
\end{itemize}

% =============================================================================
\chapter{Implémentation Backend}
% =============================================================================

\section{Endpoints requis — Comptabilité analytique}

\subsection{Écritures analytiques}

\begin{longtable}{@{}p{2cm}p{5.5cm}p{5cm}@{}}
\hline
\textbf{Méthode} & \textbf{Endpoint} & \textbf{Description} \\
\hline

GET & \texttt{/api/accounting/analytic-entries} & Liste (filtres : statut, période) \\
GET & \texttt{/api/accounting/analytic-entries/\{id\}} & Détail \\
POST & \texttt{/api/accounting/analytic-entries} & Création (accepte \texttt{clientId}) \\
PUT & \texttt{/api/accounting/analytic-entries/\{id\}} & Mise à jour \\
DELETE & \texttt{/api/accounting/analytic-entries/\{id\}} & Suppression \\
POST & \texttt{/api/accounting/analytic-entries/\{id\}/validate} & Validation \\
POST & \texttt{/api/accounting/analytic-entries/\{id\}/reject} & Rejet \\
GET & \texttt{/api/accounting/analytic-entries?since=\{iso\}} & Sync delta (phase 2) \\
\hline
\end{longtable}

\subsection{Journaux analytiques}

\begin{longtable}{@{}p{2cm}p{5.5cm}p{5cm}@{}}
\hline
\textbf{Méthode} & \textbf{Endpoint} & \textbf{Description} \\
\hline

GET & \texttt{/api/accounting/analytic-journals} & Liste \\
POST & \texttt{/api/accounting/analytic-journals} & Création \\
PUT & \texttt{/api/accounting/analytic-journals/\{id\}} & Mise à jour \\
DELETE & \texttt{/api/accounting/analytic-journals/\{id\}} & Suppression \\
\hline
\end{longtable}

\section{Contrat de création avec idempotence}

\subsection{Requête POST}

\begin{verbatim}
POST /api/accounting/analytic-entries
Headers:
  Authorization: Bearer {token}
  X-Organization-Id: {orgId}
  X-Tenant-Id: {tenantId}
  Idempotency-Key: {clientMutationId}

Body:
{
  "clientId": "ea-uuid-genere-cote-client",
  "journalId": "...",
  "dateEffet": "2026-07-08",
  "numeroPiece": "ECRIT-ANALYTIQUE-2026-0001",
  "libelleOperation": "...",
  "centreDestinationId": "...",
  "axeId": "...",
  "exerciceAnalytiqueId": "...",
  "natureChargeId": "...",
  "montant": 450000,
  "lignes": [ ... ],
  "origine": "MANUELLE"
}
\end{verbatim}

\subsection{Réponse succès}

\begin{verbatim}
HTTP 201 Created
{
  "success": true,
  "data": {
    "id": "serveur-uuid",
    "clientId": "ea-uuid-genere-cote-client",
    "version": 1,
    "updatedAt": "2026-07-08T12:00:00Z",
    ...
  }
}
\end{verbatim}

\subsection{Réponse idempotente (rejeu)}

Si \texttt{Idempotency-Key} déjà traité :

\begin{verbatim}
HTTP 200 OK
{
  "success": true,
  "data": { ... entité existante ... },
  "message": "ALREADY_PROCESSED"
}
\end{verbatim}

\section{Versionnement et conflits}

Chaque entité modifiable doit exposer :
\begin{itemize}
  \item \texttt{version} (entier incrémenté) ou \texttt{updatedAt} (ISO 8601).
  \item Le PUT/PATCH doit accepter \texttt{If-Match: \{version\}} ou \texttt{expectedVersion} dans le body.
\end{itemize}

En cas de conflit :

\begin{verbatim}
HTTP 409 Conflict
{
  "success": false,
  "code": "VERSION_CONFLICT",
  "message": "L'entité a été modifiée par un autre utilisateur",
  "data": {
    "serverVersion": { ... état serveur ... },
    "clientVersion": { ... état client ... }
  }
}
\end{verbatim}

\section{Politique de conflit recommandée (ERP comptable)}

\begin{longtable}{@{}p{4cm}p{8cm}@{}}
\hline
\textbf{Situation} & \textbf{Politique} \\
\hline

Brouillon local jamais synchronisé & Client gagne (première sync) \\
Écriture déjà validée côté serveur & Serveur gagne \\
Modification concurrente sur brouillon & Résolution manuelle (UI) \\
Suppression serveur + modification locale & Conflit, conservation locale \\
\hline
\end{longtable}

\section{Health check}

Endpoint léger pour vérifier la disponibilité API au retour en ligne :

\begin{verbatim}
GET /api/health
Response: { "status": "UP" }
\end{verbatim}

% =============================================================================
\chapter{Garanties et limites}
% =============================================================================

\section{Ce qui est garanti (niveau production)}

\begin{itemize}
  \item Survit au rechargement de page (IndexedDB persistant).
  \item Survit à une coupure réseau pendant la saisie.
  \item Sync automatique au retour en ligne.
  \item Pas de doublon si idempotence backend respectée.
  \item Opérations en attente visibles (compteur discret).
\end{itemize}

\section{Ce qui n'est pas garanti à 100\%}

\begin{itemize}
  \item Utilisateur vide le cache navigateur $\rightarrow$ perte du local non synchronisé.
  \item Navigation privée sans persistance.
  \item Multi-appareils sans sync temps réel.
  \item Quota IndexedDB dépassé (gros fichiers).
  \item Conflits non résolus manuellement.
\end{itemize}

\section{Niveaux de fiabilité}

\begin{longtable}{@{}cp{9cm}@{}}
\hline
\textbf{Niveau} & \textbf{Description} \\
\hline
A & localStorage seul, pas de backend $\rightarrow$ risque élevé \\
B & IndexedDB + outbox, un navigateur $\rightarrow$ très bon \\
C & B + backend idempotent + conflits gérés $\rightarrow$ quasi zéro perte \\
D & Absolu sur le web $\rightarrow$ impossible \\
\hline
\end{longtable}

% =============================================================================
\chapter{Plan de déploiement}
% =============================================================================

\section{Phase 1 — Fondations (actuelle)}

\begin{itemize}
  \item IndexedDB + outbox + détection réseau.
  \item OfflineProvider dans le layout.
  \item Pilote écritures analytiques.
  \item Migration localStorage $\rightarrow$ IndexedDB.
  \item Indicateur discret dans le header.
\end{itemize}

\section{Phase 2 — Lecture offline CG}

\begin{itemize}
  \item Cache des listes comptabilité générale (écritures, plan comptable, journaux).
  \item Lecture cache en repli hors ligne.
\end{itemize}

\section{Phase 3 — Écriture offline CG}

\begin{itemize}
  \item APIs backend avec idempotence.
  \item Handlers sync pour écritures comptables (brouillons d'abord).
\end{itemize}

\section{Phase 4 — PWA}

\begin{itemize}
  \item Service Worker pour shell applicatif.
  \item Ouverture de l'app sans réseau après première visite.
\end{itemize}

% =============================================================================
\chapter{Tests et validation}
% =============================================================================

\section{Scénarios de test}

\begin{enumerate}
  \item \textbf{Saisie hors ligne} : couper réseau, créer écriture, recharger page $\rightarrow$ donnée présente.
  \item \textbf{Retour en ligne} : rétablir réseau $\rightarrow$ outbox vidée, compteur à 0.
  \item \textbf{Échec API} : API down mais \texttt{onLine=true} $\rightarrow$ bascule offline après 2 échecs.
  \item \textbf{Idempotence} : rejouer même \texttt{clientMutationId} $\rightarrow$ pas de doublon.
  \item \textbf{Conflit 409} : modification concurrente $\rightarrow$ statut conflict, donnée locale conservée.
  \item \textbf{Migration} : données localStorage existantes $\rightarrow$ présentes dans IndexedDB.
\end{enumerate}

\section{Critères d'acceptation}

\begin{itemize}
  \item Aucun loader bloquant en mode offline.
  \item Polling AJAX désactivé hors ligne.
  \item Indicateur visible si opérations en attente.
  \item Aucune perte sur rechargement page en offline.
  \item Sync silencieuse au retour en ligne (toast optionnel).
\end{itemize}

% =============================================================================
\chapter{Annexes}
% =============================================================================

\section{Arborescence fichiers frontend}

\begin{verbatim}
lib/offline/
  db.ts
  network-status.ts
  outbox.ts
  sync-engine.ts
  migrate-local-storage.ts
  handlers/
    ecriture-analytique-sync.ts
hooks/
  use-network-status.ts
  use-offline-sync.ts
components/offline/
  offline-provider.tsx
  offline-status-indicator.tsx
\end{verbatim}

\section{Événements internes}

\begin{longtable}{@{}p{4cm}p{8cm}@{}}
\hline
\textbf{Événement} & \textbf{Déclencheur} \\
\hline

\texttt{network:offline} & Coupure réseau ou 2+ échecs API \\
\texttt{network:online} & Événement online + API OK \\
\texttt{sync:start} & Début flush outbox \\
\texttt{sync:complete} & Fin flush (avec count) \\
\texttt{sync:conflict} & Opération en conflit 409 \\
\hline
\end{longtable}

\section{Références}

\begin{itemize}
  \item MDN — \texttt{navigator.onLine}, IndexedDB API
  \item Dexie.js — wrapper IndexedDB
  \item Pattern Outbox — microservices / offline sync
  \item RFC — Idempotency-Key (pratique courante REST)
\end{itemize}

\end{document}
